\documentclass[11pt,a4paper]{article}
\usepackage[utf8]{inputenc}
\usepackage[T1]{fontenc}
\usepackage[polish]{babel}
\usepackage{geometry}
\geometry{margin=1.5cm}
\usepackage{longtable}
\usepackage{array}

\begin{document}

\noindent
Akademia Nauk Stosowanych w Elblągu \\
Instytut Informatyki Stosowanej im. Krzysztofa Brzeskiego

\vspace{1.5em}
\begin{center}
\textbf{\large DZIENNIK PRAKTYKI ZAWODOWEJ}
\end{center}
\vspace{1em}

\noindent Student: \dotfill\ Nr albumu: \dotfill \\
Kierunek: informatyka \hfill W zakresie: \dotfill \\
Studia: inżynierskie, stacjonarne \hfill Rok ak.: 20\dots\dots / 20\dots\dots \\
Miejsce odbywania praktyki: \dotfill \\
\dotfill \\
{\small (nazwa instytucji – zakładu pracy)} \\[1.5em]
Data rozpoczęcia praktyki: \dotfill \hfill Data zakończenia praktyki: \dotfill

\vspace{1.5em}
\noindent \textbf{WYKAZ ZAŁĄCZNIKÓW:} \\
\dotfill \\
\dotfill 

\vspace{1.5em}

% Rozciągnięcie wierszy w tabeli, aby było miejsce na odręczne wpisy (jeśli drukowane)
\renewcommand{\arraystretch}{2}

\begin{longtable}{|>{\centering\arraybackslash}p{1.5cm}|>{\centering\arraybackslash}p{2.5cm}|p{7.5cm}|>{\centering\arraybackslash}p{2.5cm}|>{\centering\arraybackslash}p{3cm}|}
\hline
\textbf{Dzień} & \textbf{Data} & \centering\textbf{Opis wykonanych prac} & \textbf{Nr. efektów uczenia się} & \textbf{Podpis osoby nadzorującej} \\ \hline
1 & 2 & \centering 3 & 4 & 5 \\ \hline
\endfirsthead

% To pojawi się na każdej kolejnej stronie po przełamaniu
\hline
\textbf{Dzień} & \textbf{Data} & \centering\textbf{Opis wykonanych prac} & \textbf{Nr. efektów uczenia się} & \textbf{Podpis osoby nadzorującej} \\ \hline
1 & 2 & \centering 3 & 4 & 5 \\ \hline
\endhead

% Stopka na dole strony, jeśli tabela przechodzi na kolejną kartkę
\hline
\endfoot

% Ostatnia stopka (na samym końcu tabeli)
\hline
\endlastfoot

% Puste wiersze przygotowane do wypełnienia (dodaj ich więcej w razie potrzeby)
 & & & & \\ \hline
 & & & & \\ \hline
 & & & & \\ \hline
 & & & & \\ \hline
 & & & & \\ \hline
 & & & & \\ \hline
 & & & & \\ \hline
 & & & & \\ \hline
 & & & & \\ \hline
 & & & & \\ \hline
 & & & & \\ \hline
 & & & & \\ \hline
 & & & & \\ \hline
 & & & & \\ \hline
 & & & & \\ \hline

\end{longtable}

\end{document}